% Critical Exponents: Power-law behaviour near P_crit
% ----------------------------------------------------------------
% Mathematical content (Theorem T-161, Appendix B):
%   Near tricritical point, order parameter
%   m_* = (r_0 / 6v)^{1/4} * epsilon^{1/4}
% i.e. m_* is proportional to epsilon^{1/4} with an amplitude that
% depends on parameters (r_0, v) of the phi^6 effective potential.
% To make the comparison between different universality classes
% meaningful, we plot y = x^beta normalised so that all curves pass
% through a common reference point at the right edge (x = x_ref).
%
% Reference point: x_ref = 0.1 (mid of the physical range, inside
% the consciousness window epsilon in (0, 1/7]).
%
% Falsification criterion (paper, Eq. below T-161):
%   beta NOT in [0.20, 0.30] at p < 0.01 --> theory falsified.
% This criterion is visualised as a band of acceptable exponents
% bounded by y = (x/0.1)^{0.20} (upper) and y = (x/0.1)^{0.30}
% (lower). Any power law whose exponent lies outside [0.20, 0.30]
% escapes this band for all x in (0, 0.1) and is therefore
% distinguishable from the UHM prediction beta = 1/4.
% ----------------------------------------------------------------

\begin{tikzpicture}
\begin{axis}[
  width=11cm, height=7cm,
  xlabel={Reduced purity $\varepsilon = P - P_{\mathrm{crit}}$},
  ylabel={Normalised order parameter
          $m/m_0 = (\varepsilon/\varepsilon_0)^{\beta}$},
  xmin=0, xmax=0.14,
  ymin=0, ymax=1.18,
  grid=both,
  grid style={gray!20},
  title={\textbf{Critical exponents of the consciousness phase
          transition}},
  title style={font=\large},
  % Legend at north-west (upper-left) with semi-transparent fill so
  % that the thick UHM curve remains visible through the legend box.
  % South-east is occupied by the 3D Ising / mean-field divergence
  % below the UHM band, and south-west is occupied by the falsification
  % boundary crossing.
  legend pos=north west,
  legend style={font=\scriptsize, fill=white, fill opacity=0.72, text opacity=1, draw=gray!50},
  every axis label/.style={font=\small},
]

% === Physical reference: epsilon_0 = 0.1 (~ 0.7 * 1/7) ===
\pgfmathsetmacro{\epsref}{0.1}

% === Falsification band: beta in [0.20, 0.30] ===
% Upper boundary (beta = 0.20) -- any beta < 0.20 goes above this.
\addplot[thin, plospurple!55, densely dashdotted,
         domain=0.001:0.14, samples=120]
  {(x/0.1)^(0.20)};
\addlegendentry{$\beta = 0.20$: UHM upper bound}

% Lower boundary (beta = 0.30) -- any beta > 0.30 goes below this.
\addplot[thin, plospurple!55, densely dashdotted,
         domain=0.001:0.14, samples=120]
  {(x/0.1)^(0.30)};
\addlegendentry{$\beta = 0.30$: UHM lower bound}

% === UHM prediction: beta = 1/4 (tricritical, T-161) ===
\addplot[ultra thick, plospurple, domain=0.001:0.14, samples=160]
  {(x/0.1)^(0.25)};
\addlegendentry{UHM: $\beta = 1/4$ (tricritical)}

% === 3D Ising: beta = 0.326 > 0.30 (falsifies UHM) ===
\addplot[thick, orange!90!black, dotted, domain=0.001:0.14, samples=160]
  {(x/0.1)^(0.326)};
\addlegendentry{3D Ising: $\beta \approx 0.326$ (falsifies UHM)}

% === Mean field phi^4: beta = 1/2 (far below UHM band) ===
\addplot[thick, gray!70!black, dashed, domain=0.001:0.14, samples=160]
  {(x/0.1)^(0.5)};
\addlegendentry{Mean field $\varphi^{4}$: $\beta = 1/2$}

% === Mark the common reference point (epsilon_0, 1) ===
\addplot[only marks, mark=*, mark size=2pt, black]
  coordinates {(0.1, 1)};
\node[font=\scriptsize, anchor=south west, inner sep=1pt]
  at (axis cs:0.1, 1.02)
  {$(\varepsilon_{0},\,1)$};

% === Goldilocks upper boundary at epsilon = 1/7 ~ 0.143 ===
% Landau expansion is valid only for small epsilon; beyond
% epsilon = 1/7 we leave the consciousness window (P > 3/7, T-124).
\addplot[thin, green!40!black, dashed] coordinates
  {(0.1429, 0) (0.1429, 1.18)};
\node[font=\scriptsize, green!40!black, anchor=south east, rotate=90]
  at (axis cs:0.1429, 0.04)
  {$\varepsilon = 1/7$: Goldilocks edge};

\end{axis}
\end{tikzpicture}
